\chapter{B-Trees and B+ Trees}
\label{ch:btree}

An index answers ``where is the data with this key?'' without reading all the
data. For data on disk the answer has been the same since 1970: the B-tree of
Bayer and McCreight \cite{bayer1972}, and its variant the B+ tree, described by
Comer as ubiquitous \cite{comer1979}. This chapter explains why it fits disks so
well, and how it stays balanced.

\section{Why a tree of pages}

A binary search tree has fanout 2: each node has two children, so finding one
of $N$ keys takes about $\log_2 N$ steps. If every step is a page read, a
billion keys need 30 reads. A B-tree makes each node a whole \emph{page} and
packs as many keys into it as fit. With fanout $F$, a lookup takes
$\log_F N$ reads.

\begin{table}[ht]
\centering
\begin{tabular}{@{}rrr@{}}
\toprule
Fanout $F$ & Height for $10^9$ keys & Height for $10^{12}$ keys \\
\midrule
2 & 30 & 40 \\
100 & 5 & 6 \\
500 & 4 & 5 \\
\bottomrule
\end{tabular}
\caption{Height is $\lceil\log_F N\rceil$. A large fanout turns a billion-key lookup into four page reads, and the top levels are almost always in the buffer pool.}
\end{table}

The upper levels are tiny (a tree of fanout 500 has 1 root and 500 children
above a quarter-million leaves) and hot, so they stay in the buffer pool. In
practice a lookup costs one or two actual disk reads.

\begin{keyidea}
A B-tree is a binary search tree redesigned for block devices: make each node a
page, make the fanout as large as the page allows, and the tree is shallow
enough that its top fits in memory.
\end{keyidea}

\section{B-tree versus B+ tree}

A \term{B-tree} stores keys and values in every node. A \term{B+ tree} stores
values only in the \term{leaves}; internal nodes hold only \term{separator keys}
that direct the search. Two consequences favor the B+ tree for databases:

\begin{itemize}
\item \textbf{Higher fanout.} Internal nodes hold no values, so more separators fit per page and the tree is shallower.
\item \textbf{Fast range scans.} Leaves are linked to their right sibling. A range query finds its first key with one descent, then walks sideways along the leaf chain without climbing back up.
\end{itemize}

\begin{figure}[ht]
\centering
\begin{tikzpicture}[font=\scriptsize]
\node[node box=3] (root) at (0,0) {\nodepart{one}20\nodepart{two}40\nodepart{three}\phantom{0}};
\node[node box=3] (l1) at (-4.2,-1.4) {\nodepart{one}5\nodepart{two}12\nodepart{three}17};
\node[node box=3] (l2) at (0,-1.4) {\nodepart{one}20\nodepart{two}31\nodepart{three}36};
\node[node box=3] (l3) at (4.2,-1.4) {\nodepart{one}40\nodepart{two}51\nodepart{three}\phantom{00}};
\draw[arrow] (root.south west) -- (l1.north);
\draw[arrow] (root.south) -- (l2.north);
\draw[arrow] (root.south east) -- (l3.north);
\draw[arrow,dashed] (l1.east) -- (l2.west);
\draw[arrow,dashed] (l2.east) -- (l3.west);
\node[below=2pt of l2] {leaves hold the keys; dashed arrows are sibling links};
\end{tikzpicture}
\caption{A B+ tree. Keys below separator 20 are in the left child; keys from 20 up to 40 in the middle; 40 and above in the right.}
\label{fig:bptree}
\end{figure}

The \term{separator rule} is the tree's central invariant: in an internal node
with separators $k_1 < k_2 < \dots < k_n$ and children $c_0, \dots, c_n$, every
key in $c_{i-1}$ is less than $k_i$, and every key in $c_i$ is at least $k_i$.

\section{Search}

Start at the root. At each internal node, binary-search the separators to pick
the child whose range contains the key, and descend. At the leaf, binary-search
for the key. A range scan $[\ell, u)$ searches for $\ell$, then follows sibling
links until a key reaches $u$.

\section{Insertion and splits}

Insert into the leaf the search reaches. If the leaf has room, done. If it is
full, \term{split} it:

\begin{enumerate}
\item Allocate a new leaf. Move the upper half of the keys into it.
\item Link it: the new leaf's sibling is the old leaf's sibling, and the old leaf's sibling becomes the new leaf.
\item Insert a separator into the parent: the first key of the new leaf, which is \emph{copied up} (it stays in the leaf).
\item If the parent is now full, split it too. An internal split \emph{pushes up} its middle key: the key moves to the grandparent and is removed from both halves.
\item If the root splits, create a new root with two children. This is the only way the tree grows taller, so every leaf is always at the same depth.
\end{enumerate}

\begin{figure}[ht]
\centering
\begin{tikzpicture}[font=\scriptsize]
\node at (-5.4,0.5) {before};
\node[node box=3] (b) at (-3.6,0) {\nodepart{one}5\nodepart{two}12\nodepart{three}17};
\node at (-1.2,0.1) {insert 14};
\draw[arrow] (-1.9,-0.2) -- (-0.5,-0.2);
\node[node box=2] (p) at (3.2,0.9) {\nodepart{one}14\nodepart{two}\phantom{0}};
\node[node box=2] (a) at (1.6,-0.5) {\nodepart{one}5\nodepart{two}12};
\node[node box=2] (c) at (4.8,-0.5) {\nodepart{one}14\nodepart{two}17};
\draw[arrow] (p.south west) -- (a.north);
\draw[arrow] (p.south) -- (c.north);
\draw[arrow,dashed] (a.east) -- (c.west);
\node at (3.2,-1.2) {14 is copied up; it stays in the right leaf};
\end{tikzpicture}
\caption{A full leaf with capacity 3 splits on insert. The new right leaf's first key becomes the separator.}
\label{fig:split}
\end{figure}

A classic bug lives in step 3: if the middle key is copied up \emph{and} left in
the \emph{left} leaf, the separator rule breaks, because the right child's keys
must be at least the separator while a copy of it sits on the left. Writing an
invariant checker before writing insert catches this at once.

\section{Deletion}

Remove the key from its leaf. If the leaf is still at least half full, done.
Otherwise:

\begin{itemize}
\item \textbf{Redistribute}: if a sibling has more than half, move a key across and update the parent's separator, which must now reflect the new boundary.
\item \textbf{Merge}: otherwise combine the leaf with a sibling, remove the separator between them from the parent, and free the empty page. The parent may now underflow, and the process repeats upward; if the root is left with one child, that child becomes the root.
\end{itemize}

Many production systems relax deletion and allow underfull nodes, merging only
when a node empties: deletes are rarer than inserts, and a tree that gets a
little sparse rarely matters.

\section{Bulk loading}

Building a tree by inserting keys one at a time leaves leaves about two thirds
full and touches pages in random order. When the keys are available up front,
\term{bulk loading} is far better: sort the keys, fill leaves completely from
left to right, link them, then build each internal level from the first keys of
the level below. One sequential pass per level, and every page is full.

\section{Iterators and pins}

A range scan returns rows one at a time to the query engine, so it is an
\term{iterator} that holds a position: the current leaf and an index inside it.
The leaf must stay pinned while the iterator points into it. When the iterator
moves to the sibling, it should unpin the current leaf \emph{before} pinning the
next one; holding both, plus a parent, can exhaust a small buffer pool and
deadlock.

\begin{ontodb}[: the B+ tree]
OntoDB's index is a B+ tree. Its key is a 24-byte \code{TripleKey}
(\cref{ch:rdf-indexing}), and the leaf stores the key itself: the triple
\emph{is} the record, so there is no separate tuple id. Leaves link to their
right sibling, the header page stores the root (or \code{INVALID\_PAGE\_ID} for
an empty tree), and a scan holds no pin and no latch between rows: it copies
the rest of a leaf's keys and lets the leaf go (design decision below).

An invariant checker (\code{CheckInvariants}) verifies sizes, key order, sibling
links, and the separator rule, and the shell can draw the tree
(\code{\textbackslash tree spo dot}). Splits and merges are multi-page
operations, so they have named crash points where a recovery test may kill the
process; \cref{ch:aries} explains why a split is redone but never undone.
Bulk load sorts ids, not strings, because the tree's order is the order of the
ids.
\end{ontodb}

\begin{decision}{Fixed slots, no key compression}
\dpart{Choice} A page header of 12 bytes follows the pageLSN: type, size, and
max size. A leaf then holds its right sibling's id and an array of 24-byte keys,
169 per 4 KiB page. An internal page holds slots of a 24-byte key and a 4-byte
child id, 144 keys and 145 children per page. Every key is stored whole.
\dpart{Alternatives} Prefix compression, which stores each key as the bytes it
does not share with its neighbour (most keys in an SPO leaf share their subject),
or a slotted layout with variable-length keys.
\dpart{Why} Fixed slots make \code{KeyAt(i)} an address computation and binary
search trivial, and a key copied out of a page is complete on its own. With
fanouts above 140, a billion triples still fit in a tree of five levels.
\dpart{Cost} Keys that share a long prefix waste most of their bytes, roughly
a third of each leaf for a subject-heavy SPO index. Leaf compression is on the
list of later work, with the question of what it does to prefix bounds.
\end{decision}

\begin{decision}{The key is the record}
\dpart{Choice} A leaf stores the 24-byte \code{TripleKey} and nothing else.
There is no value field and no record id pointing to a heap.
\dpart{Alternatives} The relational layout: a heap file of tuples, with each
index mapping a key to a record id.
\dpart{Why} A triple of ids \emph{is} its whole content. A heap would store the
same 24 bytes again and cost an extra page read per match.
\dpart{Cost} Each permutation is a complete copy of the data, so the number of
permutations is a space decision (\cref{ch:rdf-indexing}).
\end{decision}

\begin{decision}{Strict deletion}
\dpart{Choice} \code{Remove} keeps every node except the root at least half
full, by redistributing from a sibling or merging with it.
\dpart{Alternatives} Lazy deletion, as many production systems do: let nodes
become sparse and free a page only when it is empty.
\dpart{Why} The invariant checker can then assert a fill factor, and merges
become real multi-page structure changes with crash points, which the recovery
chapters need.
\dpart{Cost} More page writes on delete-heavy workloads, and the most intricate
code in the tree.
\end{decision}

\begin{decision}{The invariant checker comes before insert}
\dpart{Choice} \code{CheckInvariants} (size, order, sibling links, the
separator rule) is built and tested before \code{Insert}, and every insert and
delete test runs it.
\dpart{Alternatives} Test insert by looking up the keys afterwards.
\dpart{Why} Lookups can succeed on a tree that is already wrong; the
copied-up-key bug of this chapter does exactly that. A checker turns a
structural mistake into a failure on the operation that caused it.
\dpart{Cost} Code that never runs in production paths, and a full-tree walk in
every test.
\end{decision}

\begin{decision}{Hold the path, not parent pointers}
\dpart{Choice} Pages store no pointer to their parent. An insert or delete
keeps the pages it may have to change latched on its way down (the latch
crabbing of \cref{ch:concurrency}), so a split or merge finds its parent
already in hand.
\dpart{Alternatives} A parent pointer in every page.
\dpart{Why} A split of an internal page with 145 children would have to rewrite
the parent pointer of half of them, which is dozens of page writes for one
insert. A parent pointer is also one more thing a concurrent split can make
stale.
\dpart{Cost} The latched path is pinned: an insert into a tree whose pages are
all full pins every level, so the pool must be larger than the tree's height
plus two. The first, single-threaded version remembered page ids instead and
re-read parents; that was correct only while one thread changed the tree.
\end{decision}

\begin{decision}{Split in the middle}
\dpart{Choice} A full leaf splits into two halves, and the right half's first
key is copied up. A full internal page splits around its middle key, which
moves up.
\dpart{Alternatives} Split at the insertion point when keys arrive in
increasing order, leaving the left page full, as several engines do for
append-heavy indexes.
\dpart{Why} Triples arrive in no particular order across three permutations,
and a middle split is the simplest rule that keeps every page at least half
full.
\dpart{Cost} Inserting sorted keys one at a time leaves every leaf half full,
twice the pages a scan must read. Loading sorted data is what bulk load is for.
\end{decision}

\begin{decision}{Bulk load packs leaves full}
\dpart{Choice} \code{BulkLoad} fills each leaf to its maximum, left to right,
then builds each internal level the same way. If the last node of a level would
fall below its minimum, it and its left neighbour split their entries evenly.
\dpart{Alternatives} Leave free space in each leaf, a \emph{fill factor}, as
PostgreSQL does for its B-tree indexes (90\% by default).
\dpart{Why} An ontology is loaded and then queried. Full leaves mean the fewest
pages for every scan, and loading is one sequential pass per level.
\dpart{Cost} The first insert into a full leaf splits it at once, so a tree that
is loaded and then edited heavily ends up looking like one built by inserts.
\end{decision}

\begin{decision}{A scan holds nothing between rows}
\dpart{Choice} When the iterator reaches a leaf, it copies the leaf's remaining
keys (up to the scan's upper bound), notes the right sibling's id, and releases
the leaf. It latches the sibling only when those keys run out.
\dpart{Alternatives} Keep the current leaf pinned and read-latched for as long
as the iterator sits on it, as the first version did.
\dpart{Why} A join keeps one scan open while it opens others on other trees,
so held leaf latches were taken in both orders across trees, and Linux's
ThreadSanitizer flagged the cycle. With a writer waiting between two such
readers, that is a real deadlock, and a transaction waiting for a lock while
its scan held a latch would be one the lock manager cannot see. A scan that
holds nothing cannot take part in either.
\dpart{Cost} Up to a page of keys copied per leaf. A key inserted into a leaf
after the scan copied it is not returned, and the copied keys may include one
deleted since; the scan sees each leaf as of the moment it read it, which is
what an unlocked scan promises, and transactions add locks on top.
\end{decision}

\begin{decision}{A load arrives as one batch}
\dpart{Choice} The loader parses a file (or several) completely, interning
terms as it goes, and hands the triples to the store in one call. The indexed
store, if it is empty, sorts each permutation's keys and bulk-loads every tree;
if it already holds triples, it inserts one at a time.
\dpart{Alternatives} Insert each triple as it is parsed; or rebuild the trees
from old and new keys on every load.
\dpart{Why} Measured on LUBM(1), the 15 files in one batch load in 1.1 seconds
instead of 1.7, into 10.8 MiB instead of 16.2, and every scan reads about half
the pages, because bulk load packs the leaves that random inserts left half
full. Building three trees costs barely more than one; the time left is the
dictionary's.
\dpart{Cost} A batch holds every triple of the load in memory, and only the
first load into an empty store is bulk-loaded; later loads pay for inserts and
leave half-full leaves again.
\end{decision}

\begin{decision}{Small nodes in tests}
\dpart{Choice} The tree's maximum node sizes default to what a page holds
(169 and 144 keys), and tests pass sizes as small as 2 to grow deep trees from
a few dozen keys. A randomized test inserts and deletes against a reference set
with node sizes of 2 to 5, checking every invariant after every operation.
\dpart{Alternatives} Test only with full pages, or shrink \code{PAGE\_SIZE} for
tests.
\dpart{Why} With 169 keys per leaf, a tree needs tens of thousands of keys to
reach three levels, so bugs in internal merges and root changes would hide. The
same code runs with either size; only two numbers change.
\dpart{Cost} Small nodes make some cases more common than they are in real use
and others rarer, so the full-size tests stay as well.
\end{decision}

\begin{status}
The B+ tree is complete (plans 2.1 to 2.7): the page layout, the invariant
checker and \code{ToDot}, point search, insert with splits, the prefix
iterator, bulk load, and delete with merge and redistribution, with all four
crash points. It is safe for concurrent use by latch crabbing (plan 7.1), and
the shell's \code{\textbackslash tree} draws it on the indexed backend. Loading into
an empty indexed store goes through bulk load (plan 6.1), and LUBM(1) loads and
answers queries on a 32-frame pool (plan 6.2).
\end{status}

\begin{checkyourself}
\item OntoDB's internal page has 145 slots of 28 bytes after a 20-byte header. Where do the 28 bytes come from, and why does slot 0's key go unused?
\item A page holds 300 separators. About how many levels does a tree of $10^8$ keys need, and how many of those are likely in the buffer pool?
\item During a leaf split the separator 14 is copied up. Which leaf must it stay in, and what breaks if it stays in the other one?
\item Why does an internal split \emph{push up} its middle key, while a leaf split \emph{copies} it?
\item Why should a scan unpin its current leaf before pinning the sibling?
\item Why does OntoDB not store a parent pointer in each page? What does the choice cost once several threads change the tree?
\item OntoDB bulk-loads 170 sorted keys with a leaf maximum of 169. What leaves does it build, and why not one leaf of 169 and one of 1?
\item Bulk load fills every leaf. Why can't a loop of \code{Insert} calls do the same?
\end{checkyourself}

\section*{Further reading}
Bayer and McCreight \cite{bayer1972}; Comer \cite{comer1979}; Petrov
\cite{petrov}, chapters 2, 4, and 6.
